\subsection{扩张模与扭积的对偶性}

设 A 为环，P 与 J 为 A-模。对每个 A-模复形 C，我们在 A, X, 第 99 页，命题 12 中构造了复形的一个典范同构

\[
\mu : \operatorname{Homgr} _ {\mathrm{A}} (\mathrm{C} \otimes_ {\mathrm{A}} \mathrm{P}, \mathrm{J}) \longrightarrow \operatorname{Homgr} _ {\mathrm{A}} (\mathrm{C}, \operatorname{Hom} _ {\mathrm{A}} (\mathrm{P}, \mathrm{J})) .
\]

设 M 为\(\Lambda\) 模，(C, p) 为 M 的一个投射分解。考虑同态序列

\[
\begin{array}{r l} \operatorname{Ext} _ {\mathrm{A}} (\mathrm{M}, \operatorname{Hom} _ {\mathrm{A}} (\mathrm{P}, \mathrm{J})) & \xrightarrow {\varphi^ {- i}} \mathrm{H} (\operatorname{Homgr} _ {\mathrm{A}} (\mathrm{C}, \operatorname{Hom} _ {\mathrm{A}} (\mathrm{P}, \mathrm{J}))) \xrightarrow {\mathrm{H} (\mu) ^ {- 1}} \mathrm{H} (\operatorname{Homgr} _ {\mathrm{A}} (\mathrm{C} \otimes_ {\mathrm{A}} \mathrm{P}, \mathrm{J})) \\ & \xrightarrow {u} \operatorname{Homgr} _ {\mathrm{A}} (\mathrm{H} (\mathrm{C} \otimes_ {\mathrm{A}} \mathrm{P}), \mathrm{J}) \xrightarrow {v} \operatorname{Homgr} _ {\mathrm{A}} (\operatorname{Tor} ^ {\mathrm{A}} (\mathrm{M}, \mathrm{P}), \mathrm{J}), \end{array}
\]

其中 \(\varphi\) 是典范同构 \(\varphi(\mathrm{C},\mathrm{Hom}_{\mathrm{A}}(\mathrm{P},\mathrm{J}))\)（A, X, 第 100 页，定理 1），\(u\) 是典范同态 \(\lambda(\mathrm{C}\otimes_{\Lambda}\mathrm{P},\mathrm{J})\)（A, X, 第 82 页），而 \(v\) 由典范同构 \(\psi(\mathrm{C},\mathrm{P}):\operatorname{Tor}^{\Lambda}(\mathrm{M},\mathrm{P})\longrightarrow\mathrm{H}(\mathrm{C}\otimes_{\mathrm{A}}\mathrm{P})\) 导出。

设 \((\mathbf{C}', p')\) 为 M 的另一个投射分解。由 A, X, 第 49 页，命题 3 的推论，存在复形的一个同伦 \(\alpha: \mathbf{C}' \to \mathbf{C}\)，使 \(p \circ \alpha = p'\)。由 A, X, 第 103 页，命题 2 可知，对每个 A-模 R，有 \(\mathrm{H}(\alpha \otimes 1_{\mathrm{P}}) \circ \psi(\mathrm{C}', \mathrm{P}) = \psi(\mathrm{C}, \mathrm{P})\) 与 \(\varphi(\mathrm{C}', \mathrm{R}) \circ \mathrm{H}(\mathrm{Homgr}(\alpha, 1_{\mathrm{R}})) = \varphi(\mathrm{C}, \mathrm{R})\)。由此推出，0 次分次同态

\[
\theta (\mathrm{M}, \mathrm{P}): \operatorname{Ext} _ {\mathrm{A}} (\mathrm{M}, \operatorname{Hom} _ {\mathrm{A}} (\mathrm{P}, \mathrm{J})) \longrightarrow \operatorname{Homgr} _ {\mathrm{A}} \left(\operatorname{Tor} ^ {\mathrm{A}} (\mathrm{M}, \mathrm{P}), \mathrm{J}\right)
\]

（即上述同态序列的合成）与 M 的投射分解 (C, p) 的选取无关。按构造，它是 End \(_{A}\) (J)-线性的。

同态 \(\theta(\mathrm{M},\mathrm{P})\) 的定义可明确表述如下。设 \(p\) 为整数，\(\mathfrak{v}\) 为 \(\operatorname{Ext}_{\Lambda}^{p}(\mathrm{M},\operatorname{Hom}_{\mathrm{A}}(\mathrm{P},\mathrm{J}))\) 的一个元素，\(\tau\) 为 \(\operatorname{Tor}_{p}^{\mathrm{A}}(\mathrm{M},\mathrm{P})\) 的一个元素。利用同构 \(\varphi(\mathrm{C},\operatorname{Hom}_{\mathrm{A}}(\mathrm{P},\mathrm{J}))\)，\(\mathfrak{v}\) 由线性映射 \(u:\mathrm{C}_{p}\to\operatorname{Hom}_{\mathrm{A}}(\mathrm{P},\mathrm{J})\) 代表，且满足 \(u\circ d_{\mathrm{C}}=0\)；类似地，利用 \(\psi(\mathrm{C},\mathrm{P})\)，\(\tau\) 由元素 \(\sum c_{\mu}\otimes p_{\mu}\)（属于 \(\mathrm{C}_{p}\otimes\mathrm{P}\)，满足 \(\sum d_{\mathrm{C}}(c_{\mu})\otimes p_{\mu}=0\)）代表。于是有 \(\theta(\mathrm{M},\mathrm{P})(\mathfrak{v})(\tau)=\sum u(c_{\mu})(p_{\mu})\)。

另一方面，设 \(\nu : C \otimes_A \operatorname{Hom}_A(P, J) \longrightarrow \operatorname{Homgr}_A(\operatorname{Homgr}_A(C, P), J)\) 为把元素 \(c \otimes h\)（其中 \(c \in C_p\)，\(h \in \operatorname{Hom}_A(P, J)\)）映到同态 \(u \mapsto (-1)^p h(u(c))\) 的同态。它是 0 次分次的；若每个模 \(C_p\) 都是有限型自由模，则它是双射。不难验证它是复形的一个态射。

考虑同态序列

\[
\begin{array}{r l} \operatorname{Tor} ^ {\Lambda} (M, \operatorname{Hom} _ {\Lambda} (P, J)) & \xrightarrow {\psi} H (C \otimes_ {A} \operatorname{Hom} _ {\Lambda} (P, J)) \xrightarrow {H (\nu)} H (\operatorname{Homgr} _ {A} (\operatorname{Homgr} _ {A} (C, P), J)) \\ & \xrightarrow {w} \operatorname{Homgr} _ {\Lambda} (H (\operatorname{Homgr} _ {A} (C, P)), J) \xrightarrow {t} \operatorname{Homgr} _ {A} (\operatorname{Ext} _ {A} (M, P), J) \end{array}
\]

其中 \(\psi\) 是典范同构 \(\psi(\mathrm{C},\mathrm{Hom}_{\mathrm{A}}(\mathrm{P},\mathrm{J}))\)，\(w\) 是典范同态 \(\lambda(\mathrm{Homgr}_{\mathrm{A}}(\mathrm{C},\mathrm{P}),\mathrm{J})\)（A, X, 第 82 页），而 \(t\) 由典范同构 \(\varphi(\mathrm{C},\mathrm{P})\) 导出。与上面一样可以看出，合成同态

\[
\rho (\mathrm{M}, \mathrm{P}): \operatorname{Tor} ^ {\mathrm{A}} (\mathrm{M}, \operatorname{Hom} _ {\mathrm{A}} (\mathrm{P}, \mathrm{J})) \longrightarrow \operatorname{Homgr} _ {\mathrm{A}} \left(\operatorname{Ext} _ {\mathrm{A}} (\mathrm{M}, \mathrm{P}), \mathrm{J}\right)
\]

与分解 C 的选取无关；它是 \(\operatorname{End}_{\mathrm{A}}(\mathrm{J})\)-线性的。设 p 为整数，\(\xi \in \operatorname{Tor}_{p}^{\mathrm{A}}(\mathrm{M}, \operatorname{Hom}_{\mathrm{A}}(\mathrm{P}, \mathrm{J}))\)，\(\lambda \in \operatorname{Ext}_{\mathrm{A}}^{p}(\mathrm{M}, \mathrm{P}), \mathrm{J}\)；若 \(\xi\) 借助 \(\psi(\mathrm{C}, \operatorname{Hom}_{\Lambda}(\mathrm{P}, \mathrm{J}))\) 由元素 \(\sum c_{\mu} \otimes u_{\mu}\)（属于 \(\mathrm{C} \otimes \operatorname{Hom}_{\Lambda}(\mathrm{P}, \mathrm{J})\)，满足 \(\sum d_{\mathrm{C}}(c_{\mu}) \otimes u_{\mu} = 0\)）代表，而 \(\lambda\) 借助 \(\varphi(\mathrm{C}, \mathrm{P})\) 由同态 \(\ell : C_{p} \to P\)（满足 \(\ell \circ d_{C} = 0\)）代表，则有 \(\rho(\mathrm{M}, \mathrm{P})(\xi)(\lambda) = (-1)^{p} \sum u_{\mu} (\ell(c_{\mu}))\)。

命题 6.—— 设 A-模 J 是内射的；对每个 A-模 N，令 D(N) = \(\mathrm{Hom}_{\Lambda}(N,J)\)。

\begin{enumerate}
\def\labelenumi{\alph{enumi})}
\item
  同态 \(\theta^{i}(\mathrm{M},\mathrm{P}):\operatorname{Ext}_{\mathrm{A}}^{i}(\mathrm{M},\mathrm{D}(\mathrm{P}))\longrightarrow \mathrm{D}\big(\operatorname {Tor}_{i}^{\mathrm{A}}(\mathrm{M},\mathrm{P})\big)\) 是双射。
\item
  若环 A 是 Noether 环且 A-模 M 是有限生成的，则同态 \(\rho_{i}(\mathrm{M},\mathrm{P}):\operatorname{Tor}_{i}^{\mathrm{A}}(\mathrm{M},\mathrm{D}(\mathrm{P}))\longrightarrow \mathrm{D}\bigl (\operatorname {Ext}_{\mathrm{A}}^{i}(\mathrm{M},\mathrm{P})\bigr)\) 是双射。
\item
  按构造，只要 \(\theta(\mathrm{M},\mathrm{P})\) 是双射，同态 \(\lambda(\mathrm{C}\otimes_{\mathrm{A}}\mathrm{P},\mathrm{J})\) 便是双射的；当 J 内射时即属此种情形（A, X, 第 85 页，推论 2）。
\item
  选取分解 C，使每个模 \(C_{p}\) 都是有限型自由模（A, X, 第 53 页，命题 6）。于是同态 v 是双射，\(\lambda(\text{Homgr}_{\text{A}}(\text{C}, \text{P}), \text{J})\) 也是双射，因为 J 是内射的；因此 \(\rho(\text{M}, \text{P})\) 是双射。
\end{enumerate}

注.—— 1) 对 A-模的每个同态 \(f: \mathrm{N} \to \mathrm{N}'\)，用 \(\mathrm{D}(f): \mathrm{D}(\mathrm{N}') \to \mathrm{D}(\mathrm{N})\) 表示同态 \(\mathrm{Hom}(f, \mathbf{1}_{\mathrm{J}})\)。设 \(u: \mathrm{M} \to \mathrm{M}'\) 与 \(v: \mathrm{P} \to \mathrm{P}'\) 为 A-模的同态。选取 M 的投射分解 (C, p) 与 M' 的投射分解 (C', p')，以及复形的一个态射 \(\tilde{u}: \mathrm{C} \to \mathrm{C}'\)，使 \(p' \circ \tilde{u} = u \circ p\)（A, X, 第 49 页，命题 3）。图表

\[
\begin{array}{c c c} \operatorname{Homgr} _ {\mathrm{A}} (\mathrm{C} ^ {\prime} \otimes_ {\mathrm{A}} \mathrm{P} ^ {\prime}, \mathrm{J}) & \xrightarrow {\mu^ {\prime}} & \operatorname{Homgr} _ {\mathrm{A}} (\mathrm{C} ^ {\prime}, \operatorname{Hom} _ {\mathrm{A}} (\mathrm{P} ^ {\prime}, \mathrm{J})) \\ \operatorname{Hom} (\tilde {u} \otimes v, 1 _ {\mathrm{J}}) \Bigg \downarrow & & \Bigg \downarrow \operatorname{Hom} (\tilde {u}, \operatorname{Hom} (v, 1)) \\ \operatorname{Homgr} _ {\mathrm{A}} (\mathrm{C} \otimes_ {\mathrm{A}} \mathrm{P}, \mathrm{J}) & \xrightarrow {\mu} & \operatorname{Homgr} _ {\mathrm{A}} (\mathrm{C}, \operatorname{Hom} _ {\mathrm{A}} (\mathrm{P}, \mathrm{J})) \end{array}
\]

其中 \(\mu\) 与 \(\mu'\) 为典范同态，该图是交换的；于是由 A, X, 第 103 页，命题 2 推出一个交换图

\[
\begin{array}{c c c} \operatorname{Ext} _ {\Lambda} ^ {i} (\mathrm{M} ^ {\prime}, \mathrm{D} (\mathrm{P} ^ {\prime})) & \xrightarrow {\theta^ {i} (\mathrm{M} ^ {\prime} , \mathrm{P} ^ {\prime})} & \mathrm{D} (\operatorname{Tor} _ {i} ^ {\mathrm{A}} (\mathrm{M} ^ {\prime}, \mathrm{P} ^ {\prime})) \\ \operatorname{Ext} ^ {i} (u, \mathrm{D} (v)) \Bigg \downarrow & & \Bigg \downarrow \mathrm{D} (\operatorname{Tor} _ {i} (u, v)) \\ \operatorname{Ext} _ {\Lambda} ^ {i} (\mathrm{M}, \mathrm{D} (\mathrm{P})) & \xrightarrow {\theta^ {i} (\mathrm{M} , \mathrm{P})} & \mathrm{D} (\operatorname{Tor} _ {i} ^ {\Lambda} (\mathrm{M}, \mathrm{P})) \end{array}
\]

设 \(w: \mathrm{P}'' \to \mathrm{P}\) 为 A-模的一个同态；以类似的方式得到一个交换图

\[
\begin{array}{c c c} \operatorname{Tor} _ {i} ^ {\mathrm{A}} (\mathrm{M}, \mathrm{D} (\mathrm{P})) & \xrightarrow {\rho_ {i} (\mathrm{M} , \mathrm{P})} & \mathrm{D} (\operatorname{Ext} _ {\mathrm{A}} ^ {i} (\mathrm{M}, \mathrm{P})) \\ \operatorname{Tor} _ {i} (u, \mathrm{D} (w)) \Bigg \downarrow & & \Bigg \downarrow \mathrm{D} (\operatorname{Ext} ^ {i} (u, w)) \\ \operatorname{Tor} _ {i} ^ {\mathrm{A}} (\mathrm{M} ^ {\prime}, \mathrm{D} (\mathrm{P} ^ {\prime \prime})) & \xrightarrow {\rho_ {i} (\mathrm{M} ^ {\prime} , \mathrm{P} ^ {\prime \prime})} & \mathrm{D} (\operatorname{Ext} _ {\mathrm{A}} ^ {i} (\mathrm{M} ^ {\prime}, \mathrm{P} ^ {\prime \prime})) \end{array} .
\]

\begin{enumerate}
\def\labelenumi{\arabic{enumi})}
\setcounter{enumi}{1}
\tightlist
\item
  设
\end{enumerate}

\[
0 \to \mathrm{M} ^ {\prime} \xrightarrow {j} \mathrm{M} \xrightarrow {q} \mathrm{M} ^ {\prime \prime} \to 0\tag{ℰ}
\]

为 A-模的一个正合列。在典范自由分解上诱导的同态 \(\mathrm{L}(q):\mathrm{L}(\mathrm{M})\to\mathrm{L}(\mathrm{M}^{\prime\prime})\) 是满射，而复形 \(\operatorname{Ker}\mathrm{L}(q)\) 给出 \(M^{\prime}\) 的一个投射分解。把 A, X, 第 104 页，命题 3 应用于正合列 \(0\to\operatorname{Ker}\mathrm{L}(q)\to\mathrm{L}(\mathrm{M})\to\mathrm{L}(\mathrm{M}^{\prime\prime})\to0\)，便得到交换图

\[
\begin{array}{c c c} \operatorname{Ext} _ {\mathrm{A}} ^ {i} (\mathrm{M} ^ {\prime}, \mathrm{D} (\mathrm{P})) & \xrightarrow {\theta^ {i} (\mathrm{M} ^ {\prime} , \mathrm{P})} & \mathrm{D} (\operatorname{Tor} _ {i} ^ {\mathrm{A}} (\mathrm{M} ^ {\prime}, \mathrm{P})) \\ \delta^ {i} (\mathcal {E}, \mathrm{D} (\mathrm{P})) \Bigg \downarrow & & \Bigg \downarrow (- 1) ^ {i + 1} \mathrm{D} (\partial_ {i + 1} (\mathcal {E}, \mathrm{P})) \\ \operatorname{Ext} _ {\mathrm{A}} ^ {i + 1} (\mathrm{M} ^ {\prime \prime}, \mathrm{D} (\mathrm{P})) & \xrightarrow {\theta^ {i + 1} (\mathrm{M} ^ {\prime \prime} , \mathrm{P})} & \mathrm{D} (\operatorname{Tor} _ {i + 1} ^ {\mathrm{A}} (\mathrm{M} ^ {\prime \prime}, \mathrm{P})) \end{array}
\]

\[
\begin{array}{c c c} \operatorname{Tor} _ {i + 1} ^ {\mathrm{A}} (\mathrm{M} ^ {\prime \prime}, \mathrm{D} (\mathrm{P})) & \xrightarrow {\rho_ {i + 1} (\mathrm{M} ^ {\prime \prime} , \mathrm{P})} & \mathrm{D} (\operatorname{Ext} _ {\mathrm{A}} ^ {i + 1} (\mathrm{M} ^ {\prime \prime}, \mathrm{P})) \\ \partial_ {i + 1} (\mathcal {E}, \mathrm{D} (\mathrm{P})) \Bigg \downarrow & & \Bigg \downarrow (- 1) ^ {i + 1} \mathrm{D} (\delta^ {i} (\mathcal {E}, \mathrm{P})) \\ \operatorname{Tor} _ {i} ^ {\mathrm{A}} (\mathrm{M} ^ {\prime}, \mathrm{D} (\mathrm{P})) & \xrightarrow {\rho_ {i} (\mathrm{M} ^ {\prime} , \mathrm{P})} & \mathrm{D} (\operatorname{Ext} _ {\mathrm{A}} ^ {i} (\mathrm{M} ^ {\prime}, \mathrm{P})) \end{array}
\]

设

\[
0 \to \mathrm{P} ^ {\prime} \to \mathrm{P} \to \mathrm{P} ^ {\prime \prime} \to 0\tag{F}
\]

为 A-模的一个正合列；由于 A-模 J 是内射的，由此导出一个正合列

\[
0 \to \mathrm{D} (\mathrm{P} ^ {\prime \prime}) \to \mathrm{D} (\mathrm{P}) \to \mathrm{D} (\mathrm{P} ^ {\prime}) \to 0.\tag{\(\left(\mathcal{D}(\mathcal{F})\right)\)}
\]

把 A, X, 第 104 页，命题 3 以及第 106 页，命题 4 应用于正合列 (F) 与 (D(F))，便以类似的方式得到交换图。

\[
\begin{array}{c c c} \operatorname{Ext} _ {\Lambda} ^ {i} (\mathrm{M}, \mathrm{D} (\mathrm{P} ^ {\prime})) & \xrightarrow {\theta^ {i} (\mathrm{M} , \mathrm{P} ^ {\prime})} & \mathrm{D} (\operatorname{Tor} _ {i} ^ {\mathrm{A}} (\mathrm{M}, \mathrm{P} ^ {\prime})) \\ \delta^ {i} (\mathrm{M}, \mathrm{D} (\mathcal {F})) \Bigg \downarrow & & \Bigg \downarrow (- 1) ^ {i + 1} \mathrm{D} (\partial_ {i + 1} (\mathrm{M}, \mathcal {S})) \\ \operatorname{Ext} _ {\mathrm{A}} ^ {i + 1} (\mathrm{M}, \mathrm{D} (\mathrm{P} ^ {\prime \prime})) & \xrightarrow {\theta^ {i + 1} (\mathrm{M} , \mathrm{P} ^ {\prime \prime})} & \mathrm{D} (\operatorname{Tor} _ {i + 1} ^ {\mathrm{A}} (\mathrm{M}, \mathrm{P} ^ {\prime \prime})) \end{array}
\]

\[
\begin{array}{c c c} \operatorname{Tor} _ {i + 1} ^ {\mathrm{A}} (\mathrm{M}, \mathrm{D} (\mathrm{P} ^ {\prime})) & \xrightarrow {\rho_ {i + 1} (\mathrm{M} , \mathrm{P} ^ {\prime})} & D (\operatorname{Ext} _ {\Lambda} ^ {i + 1} (\mathrm{M}, \mathrm{P} ^ {\prime})) \\ \partial_ {i + 1} (\mathrm{M}, \mathrm{D} (\mathcal {F})) \Bigg \downarrow & & \Bigg \downarrow (\sim 1) ^ {i} \mathrm{D} (\delta^ {i} (\mathrm{M}, \mathcal {F})) \\ \operatorname{Tor} _ {i} ^ {\mathrm{A}} (\mathrm{M}, \mathrm{P} ^ {\prime \prime}) & \xrightarrow {\rho_ {i} (\mathrm{M} , \mathrm{P} ^ {\prime \prime})} & D (\operatorname{Ext} _ {\mathrm{A}} ^ {i} (\mathrm{M}, \mathrm{P} ^ {\prime \prime})) \end{array}
\]

